% intro.tex (국문 v14, 2026-10-20): 네 문단, 공백과 기여를 수사 없이 서술.
활동층은 영구동토 위에서 여름마다 녹고 겨울마다 다시 어는 지표 근처의 땅이다. 그 두께(ALT)는 계절 융해 깊이를 정하고 지반 안정성, 지표 수문, 분해에 노출되는 토양 유기탄소의 양을 좌우한다\cite{Hjort2018,Schuur2015,Biskaborn2019}. ALT는 각 해의 열 조건에 반응하므로, 연도 간 변동은 순환북극 활동층 관측망(CALM) 같은 장기 관측 프로그램이 기록하는 신호의 본질적인 부분이다\cite{Nelson1997,Streletskiy2025}. 그러나 관측은 드물고 군집되어 있다. 대부분의 측정은 소수의 집중 현장 캠페인과 장기 격자에서 나오며, 많은 지점은 한 번만 방문된다\cite{Moore2025,Veremeeva2025}.

약 1 km 해상도의 격자 ALT 제품 여러 종이 이제 북반구 영구동토 지역을 덮는다. ESA CCI Permafrost 같은 과정 기반 추정은 원격탐사 지표 온도로 구동한 지중 열 방정식을 풀고\cite{Westermann2024,Westermann2025}, 통계 · 기계학습(ML) 제품은 지점 관측을 기후 · 지형 · 토양 예측 변수에 연결한다\cite{Aalto2018,Ran2022a,Liu2024,Wei2026}. 세 가지 한계가 이 제품들의 연도별 활용을 제약한다. 첫째, 대부분의 제품은 다년 평균 예측 변수나 여러 해에 걸쳐 모은 지점-연도로 학습되어 제품 값이 특정 해의 기후에 묶여 있지 않다. 둘째, 제품마다 다른 지점 집합과 검증 방식으로 평가되어 보고된 오차(무작위 10겹 RMSE 21.6 cm\cite{Liu2024}에서 지점 제외 RMSE 60.5 cm\cite{Wei2026}까지)를 서로 비교할 수 없다. 셋째, 공간적으로 군집된 관측의 무작위 교차검증은 새 위치에서의 ML 정확도를 과대평가하므로\cite{Roberts2017,Ploton2020,MeyerPebesma2022,Wadoux2021}, 관측이 없는 지역, 즉 제품이 가장 필요한 곳에서의 성능을 가린다.

물리 기반 모델은 보완적인 경로를 제공한다. Stefan 해는 융해 깊이를 토양의 열적 성질과 수분에 의존하는 계수 하나로 그 계절의 융해 도일(TDD)에 연결한다\cite{Stefan1891,Riseborough2008,Nelson1997}. 이 해는 계절 안 융해의 진행을 재현하고 각 해의 기후에 반응하지만, 지역당 계수 하나로는 토양 유기물과 지형에 관련된 체계적 잔차가 남는다. 물리 모델이 주된 응답을 담고 학습기가 물리가 빠뜨린 부분을 보정하는 물리 기반 ML은 여러 환경 응용에서 정확도와 외삽을 개선하였고\cite{Willard2022,Read2019,Jia2021}, 알래스카에서는 Stefan 모델이 랜덤 포레스트보다 홀드아웃 지점의 정확도를 더 잘 유지하였다\cite{Gautam2025}. 그러한 결합이 기후 입력을 관측 연도에 맞출 때 Stefan 모델을 개선하는지, 라벨 없는 지역으로 전이되는지, 새 지역에 지역 라벨이 몇 개나 필요한지, 같은 셀 · 같은 해에서 공개 제품과 어떻게 비교되는지는 시험된 적이 없다.

이 연구는 이 질문들을 지점-연도 설계로 다룬다. 각 ALT 관측은 명시된 한 해의 관측 지점이고, 각 기후 입력은 그 해의 ERA5-Land 자료에서 관측일까지 누적해 계산한다\cite{munozsabater2021_era5_land}. 알래스카, 레나 델타, 캐나다 북서부의 2010--2024년 지점-연도 관측 \nncyAll{}개로 Stefan 모델, 직접 ML, 그리고 재보정 Stefan 모델의 잔차를 gradient boosting 으로 학습하는 물리 잔차 워크플로를 비교한다. 모든 방법과 공개 제품 5종은 같은 관측에 대해 같은 검증 설계로 채점하며, 여기에는 한 지역 전체의 홀드아웃과 외부 CALM 8개 지역의 시험이 포함된다. 이것이 정확도를 서로 비교할 수 있게 하는 조건이다. 이 연구는 (i) 관측 연도 기후 입력의 효과를 정량화하고, (ii) 학습기가 물리를 개선하는 조건과 물리 기준값이 전이성을 보존하는 조건을 밝히며, (iii) 새 지역에 필요한 지역 라벨 수를 제시하고, (iv) 적용 가능 영역과 예측 구간이 딸린 2011--2024년 연도별 1 km ALT 지도를 평가 프로토콜과 함께 제공한다.
